\documentclass[10pt,a4paper]{amsart}
\usepackage[margin=19mm]{geometry}
\usepackage{amsmath,amssymb,mathtools}
\usepackage[scr=rsfs,cal=euler]{mathalfa}
\usepackage{xfrac}
\usepackage[unicode,hidelinks,bookmarks=false]{hyperref}
\newcommand{\ie}{i.e.\ }
\newcommand{\cf}{cf.\ }
\newcommand{\resp}{respectively\ }
\newcommand{\Subsection}{\subsection}
\providecommand{\nobreakdash}{\nobreak\hbox{-}}
\setlength{\emergencystretch}{2em}
\multlinegap0pt
\usepackage{bm}
\usepackage{bbm}
\usepackage{dsfont}
\usepackage{xspace}
\usepackage{cancel}
\usepackage{mathtools}
\usepackage{amsthm}
\makeatletter
\@ifundefined{theorem}{\newtheorem{theorem}{Theorem}}{}
\@ifundefined{proposition}{\newtheorem{proposition}[theorem]{Proposition}}{}
\makeatother
\makeatletter
\expandafter\ifx\csname formula\endcsname\relax
\newenvironment{formula}[1]{\begin{equation}\label{eq:#1}}
                       {\end{equation}\noindent}
\fi
\expandafter\ifx\csname nota\endcsname\relax
\newenvironment{nota}[1]{\textcolor{red}
*\marginpar{\footnotesize{\textcolor{red}{#1}}}}




\fi
\makeatother
\title[The eigentheory for nonlocal cooperative-advective system an]{The eigentheory for nonlocal cooperative-advective system and its role in the study of free boundary system for directional epidemic models}
\author{Soufiane Bentout, Hoang-Hung Vo}
\date{Source: arXiv:2510.10024v1 (CC BY 4.0)}
\begin{document}
\maketitle

\noindent\emph{Source and licence.} The eigentheory for nonlocal cooperative-advective system and its role in the study of free boundary system for directional epidemic models. Version arXiv:2510.10024v1 (CC BY 4.0), distributed under the Creative Commons Attribution 4.0 International licence, as recorded in the arXiv licence metadata for this exact version and shown on the abstract page https://arxiv.org/abs/2510.10024v1. Full rights notice: licenses/arxiv-2510.10024-CC-BY-4.0.txt.\par\smallskip

\noindent\textit{Editorial scope.} Counted unit: the Proposition whose source label is \texttt{prop:monotonicity} in the article (source0.tex lines 2441--2448, source label \texttt{prop:monotonicity}), delivered together with the article's own complete original proof of it (source0.tex lines 2449--2565, one \texttt{proof} environment of 117 lines). No mathematics is written, repaired or re-derived: statement and proof are the article's own text line for line. Required context retained uncounted: eq:maineigenrvalue (lines 1036--1043). Every definition, display and label of those lines is retained unchanged. Omitted from this package and not redistributed: 1--1035, 1044--2440, 2566--5120. Nothing inside a retained excerpt is omitted or altered: every retained body is delivered exactly as the article's lines are written, so no omission receipt of any kind accompanies this package. Citations: the retained excerpts carry no citation command at all, so no bibliography entry of the article is transcribed and no reference is redistributed. Reference adaptation: none was needed and none was made; every reference inside the delivered unit resolves inside it, so no unresolved reference or citation is produced. Printed slips kept literal: no printed word, symbol, hypothesis or number of the counted statement or of its proof was altered. Layout and class: the article's own class is \texttt{amsart} with packages including geometry, babel, amsmath, amsfonts, amssymb, amsthm, relsize, setspace, nicefrac, yhmath, amscd, eucal, amsrefs, mathrsfs; the wrapper is the lane's pinned \texttt{amsart} editorial wrapper at 10pt on A4, which restates the theorem-like environments the retained text needs and adds the article's own macro definitions that the retained text needs (none), with extra wrapper packages (bm, bbm, dsfont, xspace, cancel, mathtools). The delivered numbering is the unit's own. Assets: no retained span contains a figure environment, a graphics-inclusion command, a PDF-page inclusion, a table or a TikZ picture; no image file is copied into this package, the package declares no asset dependency and no image file is redistributed with it. Licence: version arXiv:2510.10024v1 of this article is distributed under the Creative Commons Attribution 4.0 International licence, as recorded in the arXiv licence metadata for this exact version and shown on the abstract page https://arxiv.org/abs/2510.10024v1. A full rights notice is delivered at licenses/arxiv-2510.10024-CC-BY-4.0.txt. Characters: every retained excerpt is pure ASCII, so the delivered engine, XeLaTeX with its default Latin Modern text font, reports no missing character.
\begin{equation}\label{eq:maineigenrvalue}
\begin{cases}
d_1 \displaystyle\int_{-\mathcal{Z}}^{\mathcal{Z}} J_1(x-y)\psi_1(y)\,dy - d_1 \psi_1(x) - a(x) \psi_1(x) +H'(0) \psi_2(x)
+ p \psi_1^\prime(x) = \lambda \psi_1(x), & x \in [-\mathcal{Z}, \mathcal{Z}], \\[2ex]
d_2 \displaystyle\int_{-\mathcal{Z}}^{\mathcal{Z}} J_2(x-y)\psi_2(y)\,dy - d_2 \psi_2(x) - b(x) \psi_2(x) +G'(0) \psi_1(x)
+ q \psi_2^\prime(x) = \lambda \psi_2(x), & x \in [-\mathcal{Z}, \mathcal{Z}].
\end{cases}
\end{equation}

\begin{proposition}[\textbf{Monotonicity with respect to domain size}]\label{prop:monotonicity}
Assume $p,q>0$, $a,b\in L^\infty$, $H'(0),G'(0)\geq0$, and that the kernels $J_1, J_2$ satisfy assumption \textbf{(J)} so that each $\mathscr J_l$ is a positive compact integral operator on $X_l:=L^2([-\mathcal{Z}, \mathcal{Z}];\mathbb R^2)$.
Let $\lambda^\ast(\mathcal{Z})$ denote the principal eigenvalue of \eqref{eq:maineigenrvalue} on $[-\mathcal{Z}, \mathcal{Z}]$. Then
\[
\mathcal{Z}\longmapsto \lambda^\ast(\mathcal{Z})
\]
is strictly increasing and continuous on $(0,\infty)$.
\end{proposition}

\begin{proof}
Fix $0<\mathcal Z_1<\mathcal Z_2$. We compare $A_{\mathcal Z_1}(\xi)$ and $A_{\mathcal Z_2}(\xi)$ for a fixed $\xi>K$.

\medskip

View $X_{\mathcal Z_1}:=L^2([-\mathcal Z_1,\mathcal Z_1];\mathbb R^2)$ as a subspace of
$X_{\mathcal Z_2}:=L^2([-\mathcal Z_2,\mathcal Z_2];\mathbb R^2)$ by extending functions by $0$ on
$[-\mathcal Z_2,-\mathcal Z_1)\cup(\mathcal Z_1,\mathcal Z_2]$.
Under our hypotheses, for $\xi>K$ the resolvent $R_{\mathcal Z}(\xi)$ is positive on $X_{\mathcal Z}$ and $\mathscr J_{\mathcal Z}$ is a positive integral operator with kernels $J_i\ge0$.
For $f\in X_{\mathcal Z_1}$, $f\geq0$, we then have (pointwise a.e.)
\[
\big(A_{\mathcal Z_2}(\xi)f\big)_i(x)
= d_i\!\int_{-\mathcal Z_2}^{\mathcal Z_2}\! J_i(x-y)\,\big(R_{\mathcal Z_2}(\xi)f\big)_i(y)\,dy
\;\geq\; d_i\!\int_{-\mathcal Z_1}^{\mathcal Z_1}\! J_i(x-y)\,\big(R_{\mathcal Z_1}(\xi)f\big)_i(y)\,dy
= \big(A_{\mathcal Z_1}(\xi)f\big)_i(x),
\]
because: (i) $R_{\mathcal Z_2}(\xi)f\geq0$ and integrates $J_i$ over a larger set; (ii) restricting the
integral to $[-\mathcal Z_1,\mathcal Z_1]$ recovers $A_{\mathcal Z_1}(\xi)f$.
Thus, in the operator order on the cone $X_{\mathcal Z_1}^+$,
\[
0\leq A_{\mathcal Z_1}(\xi)\ \leq\ A_{\mathcal Z_2}(\xi)\quad\text{on }X_{\mathcal Z_1}.
\]

\medskip
\noindent\emph{Step 2: monotonicity of spectral radius in the cone.}
For positive compact operators on a Banach lattice, $0\leq A\leq B$ implies
$\rho(A)\leq\rho(B)$. Applying this to the restriction on $X_{\mathcal Z_1}$ yields
\[
\rho_{\mathcal Z_1}(\xi)\ \le\ \rho_{\mathcal Z_2}(\xi)\qquad\text{for every }\xi>K.
\]
Moreover, under our standard irreducibility/strong positivity hypothesis (e.g.\ kernels $J_i$ strictly positive on a set of positive measure and the resolvent kernel positive), the inequality is \emph{strict}:
\[
\rho_{\mathcal Z_1}(\xi)\ <\ \rho_{\mathcal Z_2}(\xi), \qquad \xi>K.
\]

\medskip
\noindent\emph{Step 3: solve $\rho_{\mathcal Z}(\xi)=1$ and compare $\lambda^\ast(\mathcal Z)$.}
For each fixed $\mathcal Z$, we have shown previously that
$\xi\mapsto\rho_{\mathcal Z}(\xi)$ is continuous and strictly decreasing on $(K,\infty)$, with
$\rho_{\mathcal Z}(\xi)\to+\infty$ as $\xi\downarrow K$ and $\rho_{\mathcal Z}(\xi)\to0$ as $\xi\to\infty$.
Hence there is a unique $\lambda^\ast(\mathcal Z)>K$ such that $\rho_{\mathcal Z}(\lambda^\ast(\mathcal Z))=1$.

Now fix $\mathcal Z_1<\mathcal Z_2$ and set $\lambda_1^\ast:=\lambda^\ast(\mathcal Z_1)$.
Using Step~2 at $\xi=\lambda_1^\ast$ gives
\[
1=\rho_{\mathcal Z_1}(\lambda_1^\ast)\ <\ \rho_{\mathcal Z_2}(\lambda_1^\ast).
\]
Since $\rho_{\mathcal Z_2}(\xi)$ is strictly decreasing in $\xi$, there exists a unique
$\lambda_2^\ast:=\lambda^\ast(\mathcal Z_2)$ with
$\rho_{\mathcal Z_2}(\lambda_2^\ast)=1$ and necessarily
\[
\lambda_2^\ast\ >\ \lambda_1^\ast.
\]

\medskip

Next, we prove the continuity of $\lambda^\ast(\mathcal{Z}),$
let $\mathcal Z_0>0$ be arbitrary and fix $\varepsilon>0$. We show there exists $\delta>0$ such that
\[
|\lambda^\ast(\mathcal Z)-\lambda^\ast(\mathcal Z_0)|<\varepsilon
\quad\text{whenever } |\mathcal Z-\mathcal Z_0|<\delta.
\]
Recall $\lambda^\ast(\mathcal Z)=-\xi(\mathcal Z)$ where $\xi(\mathcal Z)>K$ is the unique solution of
\[
\rho_{\mathcal Z}\big(\xi(\mathcal Z)\big)=1,\qquad
\rho_{\mathcal Z}(\xi):=\rho\!\big(\mathscr J_{\mathcal Z}(\mathscr T_{\mathcal Z}+\xi I)^{-1}\big).
\]


Since \(\xi\mapsto\rho_{\mathcal Z_0}(\xi)\) is continuous and strictly decreasing and
\(\rho_{\mathcal Z_0}(\xi(\mathcal Z_0))=1\), there exists \(\eta>0\) such that
\[
\rho_{\mathcal Z_0}\big(\xi(\mathcal Z_0)-\eta\big) > 1 \quad\text{and}\quad
\rho_{\mathcal Z_0}\big(\xi(\mathcal Z_0)+\eta\big) < 1.
\]
Because \(\rho_{\mathcal Z_0}\) is strictly decreasing these two inequalities are strict and the interval
\([\,\xi(\mathcal Z_0)-\eta,\; \xi(\mathcal Z_0)+\eta\,]\) contains no other root of \(\rho_{\mathcal Z_0}(\xi)=1\).


By the joint continuity of \((\mathcal Z,\xi)\mapsto\rho_{\mathcal Z}(\xi)\) on compact sets, there exists
\(\delta_1>0\) such that for all \(\mathcal Z\) with \(|\mathcal Z-\mathcal Z_0|<\delta_1\) and for all
\(\xi\in[\xi(\mathcal Z_0)-\eta,\xi(\mathcal Z_0)+\eta]\) we have
\[
\big|\rho_{\mathcal Z}(\xi)-\rho_{\mathcal Z_0}(\xi)\big| < \gamma,
\]
where \(0<\gamma<\min\{\rho_{\mathcal Z_0}(\xi(\mathcal Z_0)-\eta)-1,\;1-\rho_{\mathcal Z_0}(\xi(\mathcal Z_0)+\eta)\}\).
Choosing such a \(\gamma\) is possible because the two differences in the min are positive.

Hence for all \(\mathcal Z\) with \(|\mathcal Z-\mathcal Z_0|<\delta_1\) we obtain
\[
\rho_{\mathcal Z}(\xi(\mathcal Z_0)-\eta) \ge \rho_{\mathcal Z_0}(\xi(\mathcal Z_0)-\eta)-\gamma > 1,
\]
and
\[
\rho_{\mathcal Z}(\xi(\mathcal Z_0)+\eta) \le \rho_{\mathcal Z_0}(\xi(\mathcal Z_0)+\eta)+\gamma < 1.
\]

Fix any \(\mathcal Z\) with \(|\mathcal Z-\mathcal Z_0|<\delta_1\). The function \(\xi\mapsto\rho_{\mathcal Z}(\xi)\)
is continuous and strictly decreasing and we just showed it is \(>1\) at \(\xi=\xi(\mathcal Z_0)-\eta\) and \(<1\)
at \(\xi=\xi(\mathcal Z_0)+\eta\). By the intermediate value theorem there exists a root \(\xi(\mathcal Z)\in
(\xi(\mathcal Z_0)-\eta,\xi(\mathcal Z_0)+\eta)\) of \(\rho_{\mathcal Z}(\xi)=1\). By uniqueness of the root
(for fixed \(\mathcal Z\)) this must be the unique \(\xi(\mathcal Z)\). Therefore
\[
|\xi(\mathcal Z)-\xi(\mathcal Z_0)|<\eta.
\]

Recalling  \(\lambda^\ast(\mathcal Z)=-\xi(\mathcal Z)\), we obtain
\[
|\lambda^\ast(\mathcal Z)-\lambda^\ast(\mathcal Z_0)| = |\xi(\mathcal Z)-\xi(\mathcal Z_0)| < \eta.
\]
Finally choose \(\eta\) small enough so that \(\eta<\varepsilon\) and set \(\delta:=\delta_1\). Then for
\(|\mathcal Z-\mathcal Z_0|<\delta\) we have \(|\lambda^\ast(\mathcal Z)-\lambda^\ast(\mathcal Z_0)|<\varepsilon\),
proving continuity at \(\mathcal Z_0\).

Since \(\mathcal Z_0>0\) was arbitrary, \(\lambda^\ast(\mathcal Z)\) is continuous on \((0,\infty)\).

\end{proof}

\end{document}
